%% ma: L12-B07-0021
%% lop: 12
%% chuong: 2
%% bai: 7
%% dang: 12.7.9
%% loai: TN4
%% muc: NB
%% dap_an: "C"
%% nguon: "CĐ2.2.pdf giáo viên gửi 27/10/2026 (M07, Lê Minh Khôi – Nguyễn Anh Khôi, Chương II Chủ đề 2 Tọa độ của vectơ trong không gian), Câu 18"
%% kien_thuc: "Bài 7 (hình chiếu của điểm lên trục tọa độ), Bài 8 (trung điểm)"
%% trang_thai: cho-duyet
%% ly_do: "Dạng mới 12.7.9 đề xuất, chờ giáo viên duyệt"
\begin{ex}
Trong không gian $Oxyz$, cho điểm $F(6;4;1)$. Gọi $F'$ là điểm đối xứng của $F$ qua trục hoành. Tọa độ của điểm $F'$ là
\choice{$F'(-6;4;1)$}{$F'(6;4;-1)$}{\True $F'(6;-4;-1)$}{$F'(6;0;0)$}
\loigiai{Đối xứng qua $Ox$: giữ hoành độ, đổi dấu tung độ và cao độ: $F'(6;-4;-1)$.}
\end{ex}
